\section{Elliptic kernels and basis integrands}\label{app:kernel}

\paragraph{The kernel of \(P_{46}=vdI_{26}\).}
The cut of \(I_{26}\) fixes \(\rho=u/v^2\). Integrating from the regular real
diagonal determines the primitive that enters Appendix~\ref{app:pure-families},
\[
 P_{46}=\int_{x_0}^{v}\frac{\dd x}{y_\rho(x)}\,
    \mathcal K_7\big(q_+(x),q_-(x)\big),
 \qquad y_\rho(x)^2=(1+x-\rho x^2)^2-4x,
\]
\[
 x_0=\frac{4}{(1+\sqrt{1+4\sqrt\rho})^2},
 \qquad q_\pm(x)=\tfrac12\big(1-x+\rho x^2\pm y_\rho(x)\big).
\]
Here \(y_\rho(v)=d=z-\bar z\) on the upper half-plane sheet. This chart
covers \(u\leq 1\). For \(u>1\), use \(I_{26}(u,v)=u^{-2}I_{26}(1/u,v/u)\).
The table gives all 54 terms of \(\mathcal K_7=\sum_w c_w\sv_w\). Each
word denotes the single-valued iterated integral \(\sv_w=\sv_{z,w,0}\)
with endpoint \(z\), base point zero and outermost integration first.
The letter \(\bar z\) denotes the conjugate endpoint, and
\(\zeta_3=\zeta(3)\).

\begin{center}\small\begin{tabular}{@{}r l@{}}\toprule $c$ & $w$ (outermost first)\\\midrule
$-8$ & \(\mathtt{0010011}\)\quad \(\mathtt{0010100}\)\quad \(\mathtt{0110101}\)\quad \(\mathtt{1010110}\)\quad \(\mathtt{1011001}\)\\
$-4$ & \(\mathtt{0011010}\)\quad \(\mathtt{0100110}\)\quad \(\mathtt{0101001}\)\quad \(\mathtt{0110010}\)\quad \(\mathtt{1001010}\)\\
 & \(\mathtt{1010010}\)\\
$-2$ & \(\mathtt{0000110}\)\quad \(\mathtt{0001000}\)\quad \(\mathtt{0100001}\)\quad \(\mathtt{0100100}\)\quad \(\mathtt{0110000}\)\\
 & \(\mathtt{1000010}\)\\
$1$ & \(\mathtt{0000100}\)\quad \(\mathtt{0010000}\)\\
$2$ & \(\mathtt{0000101}\)\quad \(\mathtt{0001010}\)\quad \(\mathtt{0101000}\)\quad \(\mathtt{1010000}\)\\
$4$ & \(\mathtt{0010101}\)\quad \(\mathtt{0010110}\)\quad \(\mathtt{0011001}\)\quad \(\mathtt{0100010}\)\quad \(\mathtt{0100101}\)\\
 & \(\mathtt{0101100}\)\quad \(\mathtt{0110100}\)\quad \(\mathtt{1001100}\)\quad \(\mathtt{1010100}\)\\
$6$ & \(\mathtt{0010010}\)\\
$8$ & \(\mathtt{0110011}\)\quad \(\mathtt{1010011}\)\quad \(\mathtt{1011010}\)\\
\midrule
$-8$ & \(\mathtt{00}\,\bar z\,\mathtt{0110}\)\quad \(\mathtt{00}\,\bar z\,\mathtt{1001}\)\quad \(\mathtt{01}\,\bar z\,\mathtt{0011}\)\quad \(\mathtt{01}\,\bar z\,\mathtt{1100}\)\quad \(\mathtt{10}\,\bar z\,\mathtt{0011}\)\\
 & \(\mathtt{10}\,\bar z\,\mathtt{1100}\)\\
$8$ & \(\mathtt{00}\,\bar z\,\mathtt{0011}\)\quad \(\mathtt{00}\,\bar z\,\mathtt{1100}\)\quad \(\mathtt{01}\,\bar z\,\mathtt{0110}\)\quad \(\mathtt{01}\,\bar z\,\mathtt{1001}\)\quad \(\mathtt{10}\,\bar z\,\mathtt{0110}\)\\
 & \(\mathtt{10}\,\bar z\,\mathtt{1001}\)\\
\midrule
$-48\zeta_3$ & \(\mathtt{1011}\)\\
$-24\zeta_3$ & \(\mathtt{0110}\)\\
$12\zeta_3$ & \(\mathtt{0010}\)\\
$24\zeta_3$ & \(\mathtt{0011}\)\quad \(\mathtt{1001}\)\quad \(\mathtt{1010}\)\\
\bottomrule\end{tabular}\end{center}

\clearpage
\paragraph{The kernel of \(P_{49}=dI_{29}/\sqrt u\).}\label{app:i29kernel}
Holding \(v\) fixed instead gives the primitive
\[
 P_{49}=\int_{(1-\sqrt v)^2}^{u}\!
 \frac{\dd x}{\sqrt{x}\,y_v(x)}\,K_{29}\big(q_+(x),q_-(x)\big),
 \qquad y_v(x)^2=(1+x-v)^2-4x,
\]
where \(q_\pm(x)=(1+x-v\pm y_v(x))/2\) and \(y_v(u)=d\) on the
upper half-plane sheet. The lower endpoint is the regular real diagonal,
where \(P_{49}=0\). At \(v=1\) the endpoint coalesces with \(x=0\)
and the formula is read by continuity. In the table each seven-letter
word contributes \(c_w\sv_w\) to \(K_{29}\) and each four-letter word
contributes \(\zeta(3)c_w\sv_w\). The moving letters are
\(B=\bar z\) and \(C=\bar z/(\bar z-1)\).

\begin{center}\small\renewcommand{\arraystretch}{1.01}
\begin{tabular}{@{}lr@{\hspace{1.7em}}lr@{\hspace{1.7em}}lr@{\hspace{1.7em}}lr@{}}\toprule
\(w\)&\(c_w\)&\(w\)&\(c_w\)&\(w\)&\(c_w\)&\(w\)&\(c_w\)\\\midrule
\(\mathtt{0000101}\) & \(8\) & \(\mathtt{0101100}\) & \(-16\) & \(\mathtt{00B0110}\) & \(48\) & \(\mathtt{C100001}\) & \(8\) \\
\(\mathtt{0000110}\) & \(-8\) & \(\mathtt{0101101}\) & \(16\) & \(\mathtt{00B1001}\) & \(48\) & \(\mathtt{C100011}\) & \(-8\) \\
\(\mathtt{0001001}\) & \(8\) & \(\mathtt{0110000}\) & \(-16\) & \(\mathtt{00B1100}\) & \(-48\) & \(\mathtt{C100100}\) & \(-16\) \\
\(\mathtt{0001010}\) & \(-8\) & \(\mathtt{0110001}\) & \(16\) & \(\mathtt{01B0011}\) & \(24\) & \(\mathtt{C100101}\) & \(4\) \\
\(\mathtt{0001100}\) & \(16\) & \(\mathtt{0110010}\) & \(24\) & \(\mathtt{01B0110}\) & \(-24\) & \(\mathtt{C100110}\) & \(12\) \\
\(\mathtt{0001101}\) & \(-16\) & \(\mathtt{0110011}\) & \(-24\) & \(\mathtt{01B1001}\) & \(-24\) & \(\mathtt{C101001}\) & \(4\) \\
\(\mathtt{0010001}\) & \(-8\) & \(\mathtt{0110100}\) & \(-8\) & \(\mathtt{01B1100}\) & \(24\) & \(\mathtt{C101011}\) & \(-4\) \\
\(\mathtt{0010010}\) & \(-40\) & \(\mathtt{0110101}\) & \(8\) & \(\mathtt{10B0011}\) & \(24\) & \(\mathtt{C101100}\) & \(4\) \\
\(\mathtt{0010011}\) & \(48\) & \(\mathtt{1000010}\) & \(-8\) & \(\mathtt{10B0110}\) & \(-24\) & \(\mathtt{C101101}\) & \(-4\) \\
\(\mathtt{0010100}\) & \(8\) & \(\mathtt{1000100}\) & \(-8\) & \(\mathtt{10B1001}\) & \(-24\) & \(\mathtt{C110000}\) & \(8\) \\
\(\mathtt{0010101}\) & \(-4\) & \(\mathtt{1000110}\) & \(16\) & \(\mathtt{10B1100}\) & \(24\) & \(\mathtt{C110001}\) & \(-8\) \\
\(\mathtt{0010110}\) & \(-4\) & \(\mathtt{1001000}\) & \(8\) & \(\mathtt{C000011}\) & \(8\) & \(\mathtt{C110010}\) & \(-12\) \\
\(\mathtt{0011000}\) & \(24\) & \(\mathtt{1001010}\) & \(4\) & \(\mathtt{C000110}\) & \(-8\) & \(\mathtt{C110011}\) & \(12\) \\
\(\mathtt{0011001}\) & \(-36\) & \(\mathtt{1001100}\) & \(-12\) & \(\mathtt{C001001}\) & \(-16\) & \(\mathtt{C110100}\) & \(4\) \\
\(\mathtt{0011010}\) & \(12\) & \(\mathtt{1010000}\) & \(8\) & \(\mathtt{C001011}\) & \(4\) & \(\mathtt{C110101}\) & \(-4\) \\
\(\mathtt{0100001}\) & \(-8\) & \(\mathtt{1010010}\) & \(12\) & \(\mathtt{C001100}\) & \(8\) & \(\mathtt{0011}\) & \(-120\) \\
\(\mathtt{0100010}\) & \(8\) & \(\mathtt{1010011}\) & \(-24\) & \(\mathtt{C001101}\) & \(4\) & \(\mathtt{0101}\) & \(40\) \\
\(\mathtt{0100100}\) & \(24\) & \(\mathtt{1010100}\) & \(-4\) & \(\mathtt{C010010}\) & \(16\) & \(\mathtt{0110}\) & \(80\) \\
\(\mathtt{0100101}\) & \(-12\) & \(\mathtt{1010110}\) & \(8\) & \(\mathtt{C010011}\) & \(-12\) & \(\mathtt{1001}\) & \(-40\) \\
\(\mathtt{0100110}\) & \(-12\) & \(\mathtt{1011000}\) & \(-16\) & \(\mathtt{C010110}\) & \(-4\) & \(\mathtt{1010}\) & \(-40\) \\
\(\mathtt{0101000}\) & \(-8\) & \(\mathtt{1011001}\) & \(24\) & \(\mathtt{C011000}\) & \(-8\) & \(\mathtt{1011}\) & \(80\) \\
\(\mathtt{0101001}\) & \(4\) & \(\mathtt{1011010}\) & \(-8\) & \(\mathtt{C011001}\) & \(12\) & \(\mathtt{C011}\) & \(40\) \\
\(\mathtt{0101010}\) & \(4\) & \(\mathtt{00B0011}\) & \(-48\) & \(\mathtt{C011010}\) & \(-4\) & \(\mathtt{C110}\) & \(-40\) \\
\bottomrule\end{tabular}\end{center}
Gromov, Kazakov and Korchemsky's
\href{https://arxiv.org/pdf/1808.02688v1\#page=30}{two-magnon conformal-block sum
(arXiv:1808.02688v1, PDF Eqs.~(4.60)--(4.61))}
provides the native boundary data. The coefficient of \(\xi^8\) in their
reduced correlator is \(uI_{29}\), where \(\xi\) is the fishnet coupling.
The 3,141 rational equations have rank 476. The kernel reproduces 315 further nonzero Taylor coefficients, including periods, beyond the fitted orders. In the crossed limit \(w=1-z\to0\), the soft subgraph is the four-spoke wheel, giving \(I_{29}/u=-20\zeta(5)\log v+O(1)\). The crossed expansion of \(K_{29}\) gives the corresponding \(20\zeta(5)(w+\bar w)\) at linear order in the coefficient of \(\log v\). At \(z=0.35+0.25i\), the onefold and the fishnet expansion give \(I_{29}/u=142.993952659006\). Direct integration of the graph gives \(143.0050\pm0.0094\) (one standard error).

\clearpage
\paragraph{Basis integrand definitions.}
With \(s_{ij}=(x_i-x_j)^2\) and \(N=s_{13}s_{24}\), the four-loop
basis integrals used for the elliptic scalars, the finite remainders and
the Gram reduction have the normalization
\[
 I_k=N\int\prod_{j=5}^{8}\frac{\dd^4x_j}{\pi^2}\,
       \frac{n_k}{\prod_{ij\in E_k}s_{ij}},\qquad
 u=\frac{s_{12}s_{34}}{N},\quad v=\frac{s_{14}s_{23}}{N}.
\]
The last two columns together specify \(E_k\). Each pair \(ij\) denotes
one propagator \(1/s_{ij}\). In particular, legs 1 and 2 of \(I_{29}\)
both attach to vertices 7 and 8, and legs 3 and 4 both attach to 5 and 6.
Its graph numerator \(n_{29}=s_{12}s_{34}=uN\) supplies the factor \(u\)
in the conversion from the graphical function \(I_{29}/u\) drawn by Borinsky and Schnetz
in \href{https://arxiv.org/pdf/2206.10460v2\#page=30}{arXiv:2206.10460v2,
PDF Eq.~(20), p.~30} to the basis integral in Appendix~\ref{app:pure-families}.

\begin{center}\small\setlength{\tabcolsep}{3pt}
\begin{tabular}{@{}r l l l@{}}\toprule
\(k\) & \(n_k\) & External-to-internal pairs & Internal pairs\\\midrule
$26$ & $s_{13}s_{24}$ & $17\;18\;26\;28\;35\;37\;45\;46$ & $57\;58\;67\;68$ \\
$29$ & $s_{12}s_{34}$ & $17\;18\;27\;28\;35\;36\;45\;46$ & $57\;58\;67\;68$ \\
$110$ & $s_{13}s_{15}s_{24}$ & $16\;17\;18\;25\;28\;35\;37\;45\;46$ & $57\;58\;67\;68$ \\
$112$ & $s_{12}s_{13}s_{24}s_{56}$ & $16\;17\;18\;25\;26\;28\;35\;37\;45\;46$ & $57\;58\;67\;68$ \\
$113$ & $s_{13}s_{14}s_{23}s_{56}$ & $15\;16\;18\;26\;28\;35\;36\;37\;45\;47$ & $57\;58\;67\;68$ \\
$142$ & $s_{12}s_{14}s_{23}s_{56}$ & $15\;16\;18\;25\;26\;28\;36\;37\;45\;47$ & $57\;58\;67\;68$ \\
$168$ & $s_{12}s_{13}s_{14}s_{23}s_{56}^2$ & $15\;16\;17\;18\;25\;26\;28\;35\;36\;37\;45\;46$ & $57\;58\;67\;68$ \\
$214$ & $s_{14}s_{16}s_{23}s_{25}s_{34}$ & $15\;17\;18\;26\;27\;28\;35\;36\;38\;45\;46\;47$ & $56\;58\;67$ \\
$260$ & $s_{13}s_{14}s_{23}s_{24}s_{58}s_{67}$ & $16\;17\;18\;25\;27\;28\;35\;36\;38\;45\;46\;47$ & $56\;57\;68\;78$ \\
\bottomrule\end{tabular}
\end{center}
For the single-prefactor decomposition of \(I_{181}\), the auxiliary
integral \(I_{41}\) has numerator \(s_{47}s_{56}-s_{45}s_{67}-s_{46}s_{57}\)
and propagator set
\(15,16,25,27,36,37,45,46,48,57,58,67,68,78\).
The \href{https://github.com/windfolgen/svbwalkthrough/blob/main/anc/4Loop_int_basis.m}{ancillary integrand basis}
gives the full list in the same ordering.
